% =====================================================================
%  Thème 4 — CORRIGÉ — Niveau FACILE
% =====================================================================
\documentclass[11pt,a4paper]{article}
\usepackage[utf8]{inputenc}
\usepackage[T1]{fontenc}
\usepackage{lmodern}
\usepackage[french]{babel}
\usepackage{amsmath,amssymb,mathtools}
\usepackage{icomma}
\usepackage{xcolor}
\usepackage[most]{tcolorbox}
\usepackage{enumitem}
\usepackage{array}
\usepackage{booktabs}
\usepackage{fancyhdr}
\usepackage[hidelinks]{hyperref}
\usepackage[a4paper,margin=2.2cm,headheight=26pt,headsep=10pt]{geometry}

\definecolor{primary}{RGB}{223,87,99}
\definecolor{primarydark}{RGB}{150,42,55}
\definecolor{excol}{RGB}{0,135,90}

\tcbset{boxbase/.style={enhanced,breakable,boxrule=0.9pt,arc=2.5pt,
   left=3mm,right=3mm,top=2.3mm,bottom=2.3mm,
   fonttitle=\bfseries,coltitle=white,toptitle=0.6mm,bottomtitle=0.6mm}}
\newtcolorbox[auto counter]{correction}[1][]{boxbase,colback=excol!5,colframe=excol,
  title={\textsc{Corrigé}~\thetcbcounter\ifx\relax#1\relax\relax\else\textmd{ --- \itshape #1}\fi}}

\pagestyle{fancy}
\fancyhf{}
\fancyhead[L]{\small\textcolor{primarydark}{\textbf{Fiche 6} — Les limites}}
\fancyhead[R]{\small\textcolor{primarydark}{\textsc{Thème 4 — Corrigé Facile}}}
\fancyfoot[C]{\small\thepage}
\renewcommand{\headrulewidth}{0.8pt}
\renewcommand{\headrule}{\hbox to\headwidth{\color{primary}\leaders\hrule height \headrulewidth\hfill}}

\newcommand{\R}{\mathbb{R}}

\begin{document}

\begin{center}
{\Large\bfseries\color{primarydark}Thème 4 — Radicaux : facteur dominant \& conjugué}\\[2pt]
{\large\color{excol}Corrigé — Niveau Facile}
\end{center}
\vspace{6pt}

\begin{correction}[$\sqrt{x^2}=|x|$]
\begin{enumerate}[label=\alph*),leftmargin=2em,itemsep=3pt]
  \item $\sqrt{5^2}=\sqrt{25}=5=|5|$ ; \ $\sqrt{(-5)^2}=\sqrt{25}=5=|-5|$ (et \emph{non} $-5$).
  \item $\sqrt{(x-3)^2}=|x-3|=\begin{cases} x-3 & \text{si } x\geq 3,\\ -(x-3)=3-x & \text{si } x<3.\end{cases}$
  \item Pour $x>0$ : $\sqrt{9x^2}=\sqrt{9}\,\sqrt{x^2}=3|x|=3x$.
\end{enumerate}
\end{correction}

\begin{correction}[racines simples à l'infini]
\begin{enumerate}[label=\alph*),leftmargin=2em,itemsep=3pt]
  \item $\displaystyle\lim_{x\to+\infty}\sqrt{x^2+3}=\sqrt{+\infty}=+\infty$.
  \item $\displaystyle\lim_{x\to+\infty}\sqrt{x^4+x}=+\infty$ (l'intérieur $\to+\infty$).
  \item $\displaystyle\lim_{x\to+\infty}(\sqrt{x}+2)=+\infty+2=+\infty$.
\end{enumerate}
\end{correction}

\begin{correction}[écrire le conjugué]
On change le signe entre les deux termes :
\begin{enumerate}[label=\alph*),leftmargin=2em,itemsep=3pt]
  \item conjugué de $\sqrt{x}-3$ : $\ \sqrt{x}+3$.
  \item conjugué de $\sqrt{x+1}-\sqrt{x}$ : $\ \sqrt{x+1}+\sqrt{x}$.
  \item conjugué de $2-\sqrt{x}$ : $\ 2+\sqrt{x}$.
\end{enumerate}
\end{correction}

\begin{correction}[produit par le conjugué]
\begin{enumerate}[label=\alph*),leftmargin=2em,itemsep=3pt]
  \item $(\sqrt{x}-2)(\sqrt{x}+2)=(\sqrt{x})^2-2^2=x-4$.
  \item $\bigl(\sqrt{x+1}-1\bigr)\bigl(\sqrt{x+1}+1\bigr)=(x+1)-1=x$.
\end{enumerate}
\end{correction}

\begin{correction}[une première indétermination $\frac00$]
\begin{enumerate}[label=\alph*),leftmargin=2em,itemsep=3pt]
  \item En $4$ : numérateur $\sqrt{4}-2=0$, dénominateur $4-4=0$ : forme $\frac00$.
  \item $\displaystyle\frac{\sqrt{x}-2}{x-4}=\frac{(\sqrt{x}-2)(\sqrt{x}+2)}{(x-4)(\sqrt{x}+2)}
        =\frac{x-4}{(x-4)(\sqrt{x}+2)}=\frac{1}{\sqrt{x}+2}$.
  \item $\displaystyle\lim_{x\to 4}\frac{1}{\sqrt{x}+2}=\frac{1}{\sqrt4+2}=\frac{1}{4}$.
\end{enumerate}
\end{correction}

\end{document}
